\chapter{O Colapso da Linearidade, Algoritmos DLR e Simuladores Híbridos}
\label{chap:colapso_linearidade}

\section{Taxonomia e Comparativo Físico dos Algoritmos de Reconstrução}
\label{sec:taxonomia_dlr}

A evolução algorítmica da tomografia computadorizada compreende quatro gerações principais de processamento numérico de sinogramas.

A \cref{tab:taxonomia_algoritmos} detalha os algoritmos comerciais implementados pelos quatro principais fabricantes globais de tomógrafos.

\begin{table}[htbp]
  \centering
  \small
  \caption{Taxonomia dos algoritmos comerciais de reconstrução tomográfica por fabricante.}
  \label{tab:taxonomia_algoritmos}
  \begin{tabularx}{\textwidth}{lXXX}
    \toprule
    \textbf{Fabricante} & \textbf{Reconstrução Iterativa Híbrida (HIR)} & \textbf{Reconstrução Iterativa Baseada em Modelos (MBIR)} & \textbf{Reconstrução por Aprendizado Profundo (DLR)} \\
    \midrule
    \textbf{GE Healthcare} & ASiR / ASiR-V & Veo & \textbf{TrueFidelity} (treinado com FBP de dose plena) \\
    \textbf{Canon Medical} & AIDR 3D / AIDR 3D Enhanced & FIRST & \textbf{AiCE} (\emph{Advanced intelligent Clear-IQ Engine}) \\
    \textbf{Siemens Healthineers} & SAFIRE / ADMIRE & REDUCE & \textbf{Precise Image} / \textbf{Alpha Engine} (PCCT) \\
    \textbf{Philips Healthcare} & iDose4 & IMR (\emph{Iterative Model Reconstruction}) & \textbf{Precise Image} (redes convolucionais profundas) \\
    \bottomrule
  \end{tabularx}
\end{table}

\subsection{Filosofia GE Healthcare: TrueFidelity e Mapeamento Direto de FBP}
A rede neural profunda do TrueFidelity (GE Healthcare) foi concebida sob uma filosofia de preservação textural estrita. Seu treinamento supervisionado utilizou como alvo (\emph{ground truth}) imagens reconstruídas por FBP adquiridas em dose plena (alta contagem de fótons). Dessa forma, a rede aprende a remover o ruído quântico de exames de baixa dose sem alterar o formato de rampa do $NPS(f)$, preservando a granularidade visual clássica apreciada por radiologistas experientes.

\subsection{Filosofia Canon Medical: AiCE e Treinamento Ancorado em MBIR}
O algoritmo AiCE (\emph{Advanced intelligent Clear-IQ Engine}, Canon Medical) adota como alvo de treinamento imagens reconstruídas por MBIR de altíssima qualidade (algoritmo FIRST). Como o MBIR modela a óptica real do feixe e a física do ponto focal, a rede aprende a recuperar frequências espaciais elevadas em bordas ($TTF$ elevada) e atua com forte supressão de ruído em tecidos moles, exigindo calibração criteriosa para evitar a suavização excessiva de microestruturas.

\subsection{Filosofia Siemens Healthineers: Precise Image e Alpha Engine}
A Siemens Healthineers desenvolveu redes neurais que atuam de forma combinada no domínio de projeções brutas e no espaço da imagem. Em sistemas PCCT, o algoritmo \emph{Alpha Engine} processa múltiplos sinogramas espectrais simultaneamente, aplicando regularização adaptativa de gradientes que mantém a resolução espacial mesmo em níveis ultrabaixos de dose ($\text{CTDI}_{\text{vol}} < 1{,}0\text{ mGy}$).

\subsection{Filosofia Philips Healthcare: Precise Image em Espaço Híbrido}
O sistema Precise Image da Philips utiliza redes convolucionais profundas em múltiplas escalas com camadas residuais, processando dados iterativamente entre o espaço de sinogramas e o domínio espacial para eliminar artefatos de feixe endurecido e espalhamento com baixo custo computacional.

\section{A Quebra da Linearidade, Não-Estacionariedade e o Efeito Ceroso}
\label{sec:quebra_linearidade}

Em sistemas tomográficos que utilizam reconstruções por aprendizado profundo (DLR) e MBIR, a relação matemática entre as projeções brutas $\mathbf{y}$ e a imagem final $\mathbf{x}$ é estritamente não linear \cite{greffier2026, debbiche2024}:
\begin{equation}
  \mathbf{x} = \mathcal{R}_{\text{DLR}}(\mathbf{y}) \ne \mathbf{A}^\dagger \mathbf{y}
  \label{eq:nao_linear_dlr}
\end{equation}

Essa não-linearidade acarreta fenômenos físicos complexos, evidenciados na \cref{fig:dlr_non_linear}:
\begin{enumerate}
  \item \textbf{Dependência do Contraste e da Cena:} A resolução espacial da imagem deixa de ser constante, variando dinamicamente de acordo com o contraste do objeto e o nível de ruído local;
  \item \textbf{Não-Estacionariedade Espacial:} O ruído da imagem não possui propriedades estatísticas homogêneas. Em torno de bordas de alto contraste, o algoritmo preserva frequências espaciais elevadas, enquanto em regiões homogêneas de tecidos moles atua com agressiva remoção de ruído;
  \item \textbf{Efeito Ceroso (\emph{Plastic/Waxy Look}):} A excessiva concentração de energia do ruído em frequências baixas gera uma textura artificialmente lisa, que mascara lesões sutis de baixo contraste e reduz a sensibilidade de radiologistas \cite{toia2023, greffier2026};
  \item \textbf{Colapso dos Modelos Lineares Tradicionais:} O modelo analítico linear NPWE falha em prever a acurácia diagnóstica real sob reconstruções DLR, apresentando dispersão acentuada e baixa correlação ($r \approx 0{,}68$, Painel B da \cref{fig:dlr_non_linear}).
\end{enumerate}

\begin{figure}[htbp]
  \centering
  \includegraphics[width=0.98\textwidth]{figuras/fig4_dlr_non_linearity_detrending.png}
  \caption{Impacto da Não-Linearidade em DLR, Colapso de Modelos Analíticos Lineares e Metodologia de Detrending. (A) Detectabilidade $d'$ em função do nível de dose $\text{CTDI}_{\text{vol}}$ para FBP, HIR e DLR. (B) Dispersão e quebra de correlação linear do modelo NPWE ($r = 0{,}68$) versus a alta correlação do modelo DLMO ancorado na percepção de radiologistas ($r = 0{,}98$). (C) Processo de Detrending Polinomial 2D: remoção do gradiente anatômico macroscópico $P_2(x, y)$ para isolamento do ruído quântico puro $\delta I(x, y)$.}
  \label{fig:dlr_non_linear}
\end{figure}

\section{A Transição dos Simuladores: De Cilindros Homogêneos a Phantoms Antropomórficos Híbridos}
\label{sec:phantoms_hibridos}

Simuladores físicos homogêneos de acrílico ou água (como o \emph{phantom} Catphan\textregistered) foram desenvolvidos para a calibração de sistemas FBP lineares. Quando submetidos a algoritmos DLR (que foram treinados predominantemente com anatomias humanas reais), esses simuladores produzem padrões atípicos fora da distribuição de treinamento, invalidando as medições metrológicas \cite{solomon2020, racine2020}.

Para restabelecer o rigor da avaliação física, a física médica adotou a metodologia dos \emph{Phantoms} Antropomórficos Híbridos (\cref{fig:phantom_flow}).

\begin{figure}[htbp]
  \centering
  \includegraphics[width=0.95\textwidth]{figuras/flow4_phantom_hibrido_2afc.png}
  \caption{Metodologia Experimental com Phantoms Físicos Antropomórficos, Inserção Híbrida de Lesões 3D e Plataforma Psicofísica 2AFC.}
  \label{fig:phantom_flow}
\end{figure}

Essa abordagem utiliza simuladores físicos de alta fidelidade tecidual (como o \emph{phantom} FREDDIE) confeccionados por impressão 3D multimaterial para aquisição de fundos anatômicos reais ($H_0$), sobre os quais são inseridos computacionalmente modelos tridimensionais de lesões convolvidas com a PSF tridimensional do tomógrafo ($H_1$). Isso possibilita a geração de milhares de imagens com verdade de campo exata para treinamento de redes e testes psicofísicos humanos.

\section{Tratamento de Ruído em Anatomias Complexas: Detrending Polinomial 2D e Incerteza Bootstrap}
\label{sec:detrending}

Para calcular o espectro de ruído ($NPS$) no parênquima pulmonar ou no abdome de um \emph{phantom} antropomórfico, é necessário remover as variações anatômicas macroscópicas de densidade através da técnica de \emph{Detrending} Polinomial 2D (Painel C da \cref{fig:dlr_non_linear}):

Para cada sub-região de interesse $I_k(x, y)$, ajusta-se uma superfície polinomial bidimensional de 2ª ordem $P_2(x, y)$ por mínimos quadrados:
\begin{equation}
  P_2(x, y) = a_0 + a_1 x + a_2 y + a_3 x^2 + a_4 y^2 + a_5 xy
  \label{eq:polinomio_2d}
\end{equation}

A matriz de ruído puro residual $\delta I_k(x, y) = I_k(x, y) - P_2(x, y)$ é então submetida a janelamento de Hanning 2D para cálculo do $NPS$ sem vazamento espectral.

Para estimar o erro padrão e os intervalos de confiança de 95\% do índice $d'$ sem impor premissas gaussianas arbitrárias, aplica-se a técnica estatística de reamostragem Bootstrap não-paramétrica com $B = 2000$ replicações com reposição.
